\chapter{已有能力与缺失环节}

\section{食品热过程已经具备的物理描述}

下列结论直接进入后文模型，不再作为待发明的方程。

\begin{itemize}
\item 多孔介质的体积平均给出液、汽、能量的共同骨架。蒸之所以比干热更快加热内部，是因为表面冷凝释放潜热；烤之所以形成干壳，是因为蒸发吸热并降低近表面渗透率。
\item 蛋白变性、淀粉糊化、美拉德褐变的速率可用 Arrhenius 常数乘以转化率的幂，再乘以含水的开关函数。品质是反应进度，不是到达某一温度的时刻。
\item 热力安全进度沿冷点温度历史积分。冷点随射流、浮力和辐射角系数移动，不一定是几何中心。
\item 由表面温度反演内部热流是病态的，高频噪声被厚度指数放大。
\item 物理信息网络和神经算子可以加速参数化演化，它们的外推范围由训练轨迹里的执行器覆盖决定。
\item 结构因果模型里的干预，是替换结构方程的输入后再积分；相关条件期望保留了控制策略本身。
\item 模型预测控制给出有限时域的执行量，控制屏障把该执行量投影到由测量定义的安全集。
\end{itemize}

缺失的环节是这些模块之间的判据：何时后验已经够用，何时必须改变边界来买信息，信息写入哪几个参数，效应沿哪条路径改变品质，以及品质模型失效时谁来接管执行器。

\section{现有装置实际闭合了哪一段}

\begin{table}[H]
\centering
\caption{已有路径的信息止点}
\begin{tabularx}{\textwidth}{lXX}
\toprule
路径 & 已闭合的量 & 未进入后验的量 \\
\midrule
机械定时 & 无 & 全部内部状态 \\
腔温闭环 & 空气温度一点 & 冷点、含水、进度 \\
按重量调用曲线 & 申报质量与类别 & 初温、厚度、脂肪、形状 \\
穿刺探针 & 刺入点温度 & 未刺中的冷点；刺入改变局部场 \\
表面红外 & 辐射亮度 & 发射率、反射、内部梯度 \\
腔体湿度闭环 & 相对湿度设定误差 & 食品内部水分与冷凝位置 \\
图像分类后调用曲线 & 外观类别 & 中心进度；蒸汽遮挡使图像中断 \\
\bottomrule
\end{tabularx}
\end{table}

蒸与烤放在同一腔体，合并的是边界执行器，不是内部状态的估计。识别食材后调用固定曲线，曲线不随这一次的后验更新。

\section{不构成新增环节的做法}

下列做法已经覆盖常见产品行为，后文的步骤必须在此之外仍然可区分。

\begin{enumerate}
\item 按重量或图像选择时间—温度曲线。
\item 穿刺温度到达阈值后停止加热。
\item 用腔温误差调节加热器，用湿度误差调节蒸汽。
\item 用回归模型把腔温历史映射成剩余时间。
\item 漏电、超温、门锁各自独立动作，与品质估计无关。
\item 把传感器曲线映射到三维着色作为展示。
\end{enumerate}

湿度闭环最小化湿度误差。本文的信息动作最大化参数信息。二者可以驱动同一阀门，目标函数不同。展示用的温度着色不更新内部热阻，也不改变停止判据。

\section{新增环节}

\begin{table}[H]
\centering
\caption{本文加入闭环的环节}
\begin{tabularx}{\textwidth}{lXX}
\toprule
环节 & 输入与输出 & 落点 \\
\midrule
决策状态 & 全场压缩为冷点、表面、进度、热阻 & \eqref{eq:z} \\
被动估计上限 & 稳态观测不含内部热阻 & 第 \ref{ch:est} 章命题 \\
信息动作 & 相变变号脉冲后接对流脉冲 & \eqref{eq:astar} \\
路径效应 & 同一执行器的冷点项与表面项分开 & \eqref{eq:path} \\
有界更新 & 残差越界则丢弃修正项 & \eqref{eq:bnp} \\
测量隔离投影 & 安全不等式不含品质估计 & \eqref{eq:cbf} \\
\bottomrule
\end{tabularx}
\end{table}

各条线的原始文献、本文继承的方程层、以及明确不继承的层，按第 \ref{ch:lit} 章列出。后文凡是写成定理或结构方程的句子，都要能指回该章中的一条文献，或指回由该文献推出的中间引理。
